\medskip
\textbf{Partie A}
\medskip

\begin{enumerate}
\item Vérifions si les coordonnées des points vérifient ou non l'équation de $(P)$ :
	\begin{itemize}
		\item $2x_\mathrm{A} + 2y_\mathrm{A} -3z_\mathrm{A} +1= 2\times 1 + 2 \times 0 - 3 \times 1 +1= 2 + 0 - 3 +1= 0 $ :\quad A est dans le plan $(P)$;
		\item $2x_\mathrm{B} + 2y_\mathrm{B} -3z_\mathrm{B} +1= 2\times 2 + 2 \times (-1) - 3 \times 1+1 = 4 -2 - 3 +1= 0$ :\quad B est dans le plan $(P)$;
		\item $2x_\mathrm{C} + 2y_\mathrm{C} -3z_\mathrm{C} +1= 2\times (-4) + 2 \times (-6) - 3 \times 5+1 = -8 -12 - 15 +1= -34 \neq 0$ :\quad C n'est pas dans le plan $(P)$.
	\end{itemize}

\item On a:\quad $ \vect{\mathrm{CC'}}\begin{pmatrix} 0-(-4)\\-2-(-6)\\-1-5 \end{pmatrix} = \begin{pmatrix}4\\4\\-6\end{pmatrix} $, d'après l'équation que l'on a du plan $ (P) $, $ \vect{n}\begin{pmatrix} 2\\2\\-3 \end{pmatrix} $ est un vecteur normal au plan $ (P) $.

Comme on a manifestement $ \vect{\mathrm{CC'}} = 2\vect{n} $, ces vecteurs sont colinéaires, et donc la droite (CC') est orthogonale au plan $ (P) $.

De plus :
\quad
	$2x_\mathrm{C'} + 2y_\mathrm{C'} -3z_\mathrm{C'} + 1 = 2\times 0 + 2 \times (-2) - 3 \times (-1) + 1 = 0 - 4 + 3 + 1 = 0$ :\quad C' $\in(P)$.

Finalement, C' est un point du plan $ (P) $ tel que (CC') est orthogonale à $(P)$ : cela confirme que C' est le projeté orthogonal de C sur $(P)$.

	\item La droite (AB) passe par A$(1~;~0~;~1)$ et est dirigée par $\vect{\mathrm{AB}} \begin{pmatrix}1\\-1\\0\end{pmatrix}$, donc une représentation paramétrique de (AB) est :
	\quad
	$\begin{cases}x = x_\mathrm{A} + x_{\mathrm{\vect{\mathrm{AB}}}}t\\y = y_\mathrm{A} + y_{\mathrm{\vect{\mathrm{AB}}}}t\\z = z _\mathrm{A} + z _{\mathrm{\vect{\mathrm{AB}}}}t \end{cases}$ avec $t \in\R$,\quad soit ici :\quad $\begin{cases} x = 1 + t\\y= -t\\z = 1	\end{cases}$ avec $t \in \R$.

	\item Si H est un point de (AB), cela signifie qu'il existe un réel $t_0$ tel que H est le point de paramètre $t_0$ sur (AB).

On a : (AB) et (CH) seront orthogonales (et donc nécessairement perpendiculaires, car elles ont H en commun) si et seulement si les vecteurs $\vect{\mathrm{AB}}$ et $\vect{\mathrm{CH}}$ sont orthogonaux.

On a les coordonnées suivantes : \quad $\vect{\mathrm{AB}} \begin{pmatrix}1\\-1\\0\end{pmatrix}$ et $\vect{\mathrm{CH}}\begin{pmatrix}(1 + t_0) - (-4)\\(-t_0) - (-6) \\1 - 5\end{pmatrix} =\begin{pmatrix}5 + t_0 \\ 6 - t_0 \\ -4\end{pmatrix}$. Comme le repère est orthonormé, on a :\quad $\vect{\mathrm{AB}}\cdot\vect{\mathrm{CH}} = 1\times (5 + t_0) + (-1) \times (6 - t_0) + 0\times (-4) = 5 + t_0 - 6 + t_0 = 2t_0 - 1$

On a donc :\quad (AB) et (CH) sont orthogonales $\begin{aligned}[t] &\iff \vect{\mathrm{AB}}\cdot \vect{\mathrm{CH}} = 0\\
	&\iff 2t_0 - 1 = 0\\
	&\iff t_0 = \dfrac{1}{2}\end{aligned}$

L'unique pour H de la droite (AB) pour lequel (AB) et CH) sont orthogonales est donc le point de paramètre $t_0 = \dfrac{1}{2}=0,5$ dans la représentation que nous avons donnée à la question précédente. Ses coordonnées sont donc : H$(1,5 ~;~ -0,5 ~;~ 1)$.

\emph{Remarques :} H est le milieu de [AB].

Si vous avez une représentation paramétrique différente de celle présentée dans ce corrigé, vous aurez aussi une valeur différente pour $t_0$.

Par exemple, si vous avez choisi B comme point de référence et $\vect{\mathrm{AB}}$ comme vecteur directeur pour construire votre représentation paramétrique, vous arriverez à $t_0 = -0,5$, mais dans votre représentation paramétrique, cela vous conduira aux mêmes coordonnées pour le point H.
\end{enumerate}

\medskip

\textbf{Partie B}

\medskip
Les coordonnées admises pour le vecteur $\vect{\mathrm{HC}}$ sont cohérentes avec les coordonnées du point H à la question précédente.

\begin{enumerate}
	\item On est dans un repère orthonormé, donc :

	$\left\|\vect{\mathrm{HC}}\right\| = \sqrt{\left(-\dfrac{11}{2}\right)^2+\left(-\dfrac{11}{2}\right)^2+4^2 } = \sqrt{\dfrac{121}{4}+\dfrac{121}{4}+16} = \sqrt{\dfrac{153}{2}} = \sqrt{76,5}$.

	La valeur exacte de $\left\|\vect{\mathrm{HC}}\right\| $ est donc : $\sqrt{\dfrac{153}{2}} = \sqrt{76,5}$.

	\item Le point H, en tant que point de (AB) tel que (CH) est orthogonale à (AB), est donc le pied de la hauteur issue de C dans le triangle ABC.

Pour calculer la surface du triangle, on va donc utiliser la formule :\quad$\mathcal{A} = \dfrac{\mathrm{base}\times \mathrm{hauteur}}{2}$.

Ici, le plus simple sera donc de choisir [AB] comme base, de longueur AB et donc la hauteur correspondante est $\mathrm{CH} = \left\|\vect{\mathrm{CH}}\right\| $.

On a déjà calculé CH à la question précédente, donc :

$\mathrm{AB} = \sqrt{1^2 + (-1)^2 + 0} = \sqrt{2}$.

L'aire de ABC est donc : \quad$ S = \dfrac{\mathrm{AB}\times \mathrm{CH}}{2} = \dfrac{\sqrt{2} \times \sqrt{\dfrac{153}{2}}}{2} = \dfrac{\sqrt{153}}{2}$.
\end{enumerate}

\medskip

\textbf{Partie C}

\medskip

\begin{enumerate}
	\item \begin{itemize}[left=0mm]
		\item On a établi au début de l'exercice que A et B appartiennent à $(P)$, donc toute la droite (AB) est incluse dans $(P)$, et donc notamment H appartient à $(P)$ aussi.

		\item H et C$'$ sont donc deux points de $(P)$, et on sait que (CC$'$) est orthogonale à $(P)$.
	\end{itemize}

	(CC') étant orthogonale à $(P)$, elle orthogonale à toute droite de $(P)$, dont la droite (C'H), ces droites sont donc perpendiculaires, car elles se coupent évidemment en C'.

	Le triangle CHC' est donc un triangle rectangle en C'.

	Dans un triangle rectangle, le cosinus d'un angle aigu est donné par le quotient de la longueur du côté adjacent à l'angle (ici C$'$H) par la longueur de l'hypoténuse du triangle (ici CH).

	Ici, on a donc : \quad $\cos(\alpha ) = \dfrac{\mathrm{C'H}}{\mathrm{CH}} = \dfrac{\sqrt{\dfrac{17}{2}}}{\sqrt{\dfrac{153}{2}} } = \sqrt{\dfrac{17}{153}} = \sqrt{\dfrac{17\times 1}{17 \times 9}} = \sqrt{\dfrac{1}{9}} = \dfrac13$.

	\item \begin{enumerate}
		\item \emph{Méthode $1$ :} le plus rapide, ici, serait de calculer le produit scalaire des vecteur directeurs des deux droites, et de prouver qu'il est nul.

		\smallskip

Mais comme on a fait quelque chose de similaire dans ce corrigé, on va explorer une voie différente :


	\emph{Méthode $2$ :}
	\begin{itemize}[left=0mm]
		\item On sait que (CH) est orthogonale à (AB), d'après la définition du point H à la question \textbf{A 4.};
		\item on sait que (CC$'$) est orthogonale à (AB), car (CC$'$) étant orthogonale à $(P)$, elle est orthogonale à toute droite incluse dans $(P)$, notamment (AB);
		\item on sait que les droites (CH) et (CC$'$) sont sécantes (en C), car les points C$'$ et H sont les intersections de ces deux droites avec $(P)$, et sont séparés par une distance non nulle.
	\end{itemize}

Ainsi, on vient de démontrer que C, C$'$ et H définissent un plan, et que deux droites sécantes du plan (CC'H) sont orthogonales à (AB), donc que (AB) est orthogonale au plan (CC$'$H) et donc à toutes les droites de ce plan là, notamment à la droite (C$'$H).

		\emph{Remarque :} le plan (CC$'$H) est ce que l'on appelle le \textbf{plan médiateur} du segment [AB] : le plan qui passe par le milieu (H) du segment et lui est perpendiculaire.

		\item On applique à nouveau la formule de calcul de l'aire d'un triangle.

Ici, on va prendre [AB] comme base et donc, d'après la question précédente, [C$'$H] est la hauteur correspondante.

On a donc :\quad $S' = \dfrac{\mathrm{AB}\times \mathrm{C'H}}{2} = \dfrac{\sqrt{2} \times \sqrt{\dfrac{17}{2}}}{2} = \dfrac{\sqrt{17}}{2}$.

		\item On a donc :\quad $\cos(\alpha ) = \dfrac{\mathrm{C'H}}{\mathrm{CH}} = \dfrac{\mathrm{AB} \times \mathrm{C'H}}{\mathrm{AB} \times \mathrm{CH}}= \dfrac{\dfrac{\mathrm{AB} \times \mathrm{C'H}}{2}}{\dfrac{\mathrm{AB} \times \mathrm{CH}}{2}} = \dfrac{S'}{S}$.

Voilà une relation possible.

On peut aussi écrire :\quad $S' = \cos(\alpha)S$\quad par exemple.
	\end{enumerate}
\end{enumerate}
